\section{Model geometry and approach outline}\label{sec:geometry}
To model diffusion of insulin from capillary to neighboring cells,
we introduce a geometry of two concentric cylinders,
with a constant ligand concentration imposed at the inner cylinder of radius $R_1$ and absorbing insulin receptors distributed on the outer cylinder of radius $R_2$ (figure~\ref{fig:capillary_cells}).
We specify the associated boundary conditions and replace the heterogeneous receptor distribution with an effective homogenized boundary condition,
enabling analytical treatment of the diffusion problem.
We then incorporate receptor binding and unbinding kinetics and, finally, validate the resulting analytical predictions by comparison with stochastic random-walk simulations.

\begin{figure}[H]
    \centering
    % Placeholder for Figure 1: 3D representation of the capillary and surrounding cells.
    \fbox{\parbox[c][6cm][c]{0.8\textwidth}{\centering \textit{Placeholder for Figure 1: 3D representation of the capillary and surrounding cells.}}}
    \caption{\textbf{3D representation of the capillary and surrounding cells.} The orange disks denote receptors uniformly distributed along the cell membrane.}
    \label{fig:capillary_cells}
\end{figure}

\begin{table}[H]
    \centering
    \caption{\textbf{System variables and parameters used in the model.}}
    \label{tab:parameters}
    \begin{tabular}{@{}lll@{}}
        \toprule
        \textbf{Variable} & \textbf{Value} & \textbf{Description} \\
        \midrule
        \multicolumn{3}{@{}l}{\textit{Total body variables}} \\
        $C_{\mathrm{insulin, blood}}$ & $66\,\mathrm{pM}$ & Blood insulin concentration \\
        $N_{\mathrm{cells, body}}$ & $2 \times 10^{12}$ & Total number of cells in the body \\
        $V_{\mathrm{blood}}$ & $5\,\mathrm{L}$ & Total blood volume \\
        $N_{\mathrm{insulin}}/N_{\mathrm{cells, body}}$ & $\sim 100$ & Insulin molecules per cell \\
        \multicolumn{3}{@{}l}{\textit{General system variables}} \\
        $R_{\mathrm{cell}}$ & $5\,\mu\mathrm{m}$ & Cell radius \\
        $R_{\mathrm{outer\,ring}}$ & $7.5\,\mu\mathrm{m}$ & Cell ring radius \\
        $R_{\mathrm{capillary}}$ & $5\,\mu\mathrm{m}$ & Capillary radius \\
        $L_{\mathrm{ring}}$ & $10\,\mu\mathrm{m}$ & Cylinder length \\
        $A_{\mathrm{ring}}$ & $471\,\mu\mathrm{m}^2$ & Ring surface area \\
        $N_{\mathrm{cells\,on\,ring}}$ & $\approx 4$ & Cells per ring \\
        $R_{\mathrm{receptor}}$ & $5\,\mathrm{nm}$ & Receptor radius \\
        $N_{\mathrm{receptors\,per\,cell}}$ & $10^{4}$ & Receptors per cell \\
        $\sigma_r$ & $84.9\,\mu\mathrm{m}^{-2}$ & Receptor surface density \\
        $D_{\mathrm{insulin}}$ & $125\,\mu\mathrm{m}^{2}/\mathrm{s}$ & Diffusion coefficient \\
        \multicolumn{3}{@{}l}{\textit{Insulin--receptor kinetics}} \\
        $K_{d, \mathrm{insulin}}$ & $2.7\,\mathrm{nM}$ & Dissociation constant \\
        $k_a$ & $3 \times 10^{6}\,\mathrm{M}^{-1}\mathrm{s}^{-1}$ & Association (binding) rate [CITE] \\
        $k_b$ & $4 \times 10^{-3}\,\mathrm{s}^{-1}$ & Dissociation (unbinding) rate [CITE] \\
        $k_c$ & $4 \times 10^{-3}\,\mathrm{s}^{-1}$ & Internalization rate [CITE] \\
        \bottomrule
    \end{tabular}
\end{table}
